% ScholarWeave LaTeX export — "article" doc type.
% See document.tex for the full explanation of how this file is used
% (--include-in-header, available substitutions/hooks/markers — including
% SWTOKFIGURECOUNTER and SWTOKNEWPAGEHEADING below, driven by export
% checkboxes rather than user-edited — and title block vs. title page). A
% distinct file per doc type (mirroring the DOCX/ODT templates) so article
% and document can diverge further later without affecting each other —
% today they differ only in font and link color, matching
% article.docx/article.odt vs. document.docx/document.odt.

% SWEXTRACLASSOPTIONS:

% ── Customize this template ────────────────────────────────────────────
% Copy this file (e.g. into your own Export Templates folder) and edit any
% of the lines below to get your own fonts/margins/link color/spacing.
%
% Noto Serif here is a VARIABLE font: one file with a continuous `wght` axis
% (100–900) and no separate Bold file. LuaLaTeX exposes that axis through
% fontspec's Weight= key, so BoldFeatures={Weight=700} below is the font's
% REAL 700 instance — not a synthetic fake-bold. Weight=400, the axis
% default, is normal text. (fontspec under XeLaTeX can't address variable
% axes, so the \else branch keeps the old synthetic FakeBold, used only when
% a .tex is compiled by hand with xelatex.) Width=/Slant=/RawAxis= work the
% same way if you ever need them.
\ifluatex
  \setmainfont{Noto Serif}[BoldFont={Noto Serif},BoldFeatures={Weight=700}]
\else
  \setmainfont{Noto Serif}[BoldFont={Noto Serif},BoldFeatures={FakeBold=3.5}]
\fi
\usepackage[letterpaper,margin=1.3in]{geometry}  % page size + margins (try a4paper, or a per-side margin like "top=1in,bottom=1in,left=1.25in,right=1in")
\definecolor{swlinkcolor}{HTML}{467886}        % link/citation color, as a hex RRGGBB (no '#') — matches article.docx/.odt's Hyperlink style
\usepackage{setspace}\setstretch{1.35}         % paragraph line spacing — a setspace MULTIPLE (1 = single); this is NOT the same number as the DOCX/ODT 1.15 line-spacing multiple (they measure from a different baseline reference), so tune it by eye.
\SetSinglespace{1.35}                          % FOOTNOTE/FLOAT line spacing — setspace forces footnotes and floats to SINGLE spacing regardless of \setstretch, so they look cramped next to the body; set this to match \setstretch (or lower). Affects footnotes & floats only, not body text.
\raggedbottom                                  % let pages end wherever content naturally does — see document.tex
                                                % for why (stretchy spacing otherwise fills to an even bottom
                                                % margin, at the cost of inconsistent-looking gaps). Switch to
                                                % \flushbottom for even bottom margins instead.
% ────────────────────────────────────────────────────────────────────────

% Title-block logo hook — empty by default. Redefine to add a journal/
% institution logo above the title, e.g.:
%   \renewcommand{\swTitleLogo}{\includegraphics[width=3cm]{/path/to/logo.png}\\[1em]}
\newcommand{\swTitleLogo}{}

% ── Arabic / non-Latin RTL scripts (wrapper-free) ───────────────────────
% See document.tex for the full rationale. In short: LuaLaTeX + babel
% `bidi=basic` gives correct Arabic letter-joining AND right-to-left
% ordering (XeLaTeX cannot — its macro-based alternatives crash with "TeX
% capacity exceeded" at the first footnote), and `onchar=ids fonts`
% auto-detects any Arabic-script run (including Arabic-derived "Ajami") with
% NO wrapper, so inline Arabic just works. The Arabic font is Scheherazade
% New at 120% (Arabic looks smaller than Latin at the same nominal size),
% with Noto Naskh Arabic as the fallback; Chinese is provided too (font
% selection only — CJK needs no bidi). The nested \IfFontExistsTF guards
% let a missing font degrade instead of breaking the export.
%
% Right indent for Arabic verse lines (babel maps LaTeX's "left" indents to
% the RTL right side) — the verse default is nearly flush. Raise it for a
% deeper indent.
\newlength{\swArabicVerseIndent}\setlength{\swArabicVerseIndent}{5em}
\newcommand{\swArabicVerseSetup}{\setlength{\leftmargini}{\swArabicVerseIndent}}

\ifluatex
  \makeatletter
  \@ifpackageloaded{babel}{}{\usepackage[bidi=basic,shorthands=off]{babel}}
  \makeatother
  \babelprovide[import,onchar=ids fonts]{arabic}
  \IfFontExistsTF{Scheherazade New}{%
    \babelfont[arabic]{rm}[Scale=1.2]{Scheherazade New}%
  }{\IfFontExistsTF{Noto Naskh Arabic}{\babelfont[arabic]{rm}[Scale=1.2]{Noto Naskh Arabic}}{}}%
  \babelprovide[import,onchar=ids fonts]{chinese}
  \IfFontExistsTF{Noto Serif CJK SC}{%
    \babelfont[chinese]{rm}{Noto Serif CJK SC}%
  }{\IfFontExistsTF{Songti SC}{\babelfont[chinese]{rm}{Songti SC}}{}}%
  \newcommand{\arabicfont}{\selectlanguage{arabic}\swArabicVerseSetup}
\else
  % XeLaTeX fallback (manual .tex compile only): no paragraph-level bidi, so
  % mixed inline runs come out reversed. LuaLaTeX is strongly preferred.
  \newfontfamily\swArabicFont[Script=Arabic]{Scheherazade New}
  \newcommand{\arabicfont}{\swArabicFont\swArabicVerseSetup}
\fi

% Bullet-list levels alternate bullet/dash instead of LaTeX's own default
% bullet/dash/asterisk/centered-dot cascade.
\usepackage{enumitem}
\setlist[itemize,1]{label=\textbullet}
\setlist[itemize,2]{label=\textendash}
\setlist[itemize,3]{label=\textbullet}
\setlist[itemize,4]{label=\textendash}

% Heading 1 (\section) / Heading 2 (\subsection) styles — edit these to
% change appearance (bold/size/spacing/etc.); deeper heading levels use
% LaTeX's own defaults. No number prefix by default.
%
% \titlespacing's before/after values are FIXED lengths (no "plus/minus"
% stretch component) deliberately — see document.tex for why: the class's
% own default heading spacing is stretchy, so a heading's gap from the text
% above it can end up much larger on one page than another with no visible
% cause, whenever the page-breaking algorithm stretches it to fill a page.
\usepackage{titlesec}
\titleformat{\section}{\normalfont\Large\bfseries}{}{0pt}{SWTOKNEWPAGEHEADING}
\titlespacing*{\section}{0pt}{3.5ex}{2.3ex}
\titleformat{\subsection}{\normalfont\large\bfseries}{}{0pt}{}
\titlespacing*{\subsection}{0pt}{3.25ex}{1.5ex}
% Headings show no number (above), and the TOC entries below match — but
% \thesection still increments normally (--metadata numbersections=true,
% set in DocumentCompiler.py): SWTOKFIGURECOUNTER below needs a real,
% incrementing section counter to scope figures by, which secnumdepth's
% usual "no visible number" approach silently prevents (it stops the
% counter, not just the display).
\usepackage{titletoc}
\titlecontents{section}[0pt]{}{}{}{\hfill\thecontentspage}
\titlecontents{subsection}[1.5em]{}{}{}{\hfill\thecontentspage}
SWTOKFIGURECOUNTER

\usepackage{fancyhdr}
\pagestyle{fancy}
\fancyhf{}
\fancyhead[LE]{SWTOKAUTHOR}
\fancyhead[RO]{SWTOKSHORTTITLE}
\fancyfoot[C]{\thepage}
\renewcommand{\headrulewidth}{0.4pt}
% Page 1 (the title block) uses LaTeX's "plain" pagestyle — redefine it to
% keep the centered page-number footer but drop the running header. See
% document.tex for the full explanation.
\fancypagestyle{plain}{%
  \fancyhf{}
  \fancyfoot[C]{\thepage}
  \renewcommand{\headrulewidth}{0pt}
}
% See document.tex for why \hypersetup{pdftitle=...} isn't set here, and for
% how colorlinks/linkcolor above still reaches pandoc's own hypersetup call.
